% SPDX-License-Identifier: CC-BY-4.0
\documentclass[11pt,letterpaper]{article}

\usepackage[T1]{fontenc}
\usepackage{amsmath,amssymb}
\usepackage{tikz}
\usepackage{pgfplots}
\pgfplotsset{compat=1.18}
\usepackage{algorithm}
\usepackage{algpseudocode}
\usepackage{booktabs}
\usepackage{enumitem}
\usepackage[hidelinks]{hyperref}

\newcommand{\N}{\mathbb{N}}
\newcommand{\Z}{\mathbb{Z}}
\newcommand{\Q}{\mathbb{Q}}
\newcommand{\R}{\mathbb{R}}
\newcommand{\C}{\mathbb{C}}
\newcommand{\Map}[3]{#1 \colon #2 \to #3}
\newcommand{\Mod}[1]{#1\text{-}\mathrm{Mod}}
\newcommand{\GroupHom}{\mathrm{Hom}_{\mathbf{Grp}}}
\newcommand{\RingHom}{\mathrm{Hom}_{\mathbf{Ring}}}

\newtheorem{theorem}{Theorem}
\newtheorem{lemma}[theorem]{Lemma}
\newtheorem{corollary}[theorem]{Corollary}
\newtheorem{definition}[theorem]{Definition}
\newtheorem{remark}[theorem]{Remark}
\newenvironment{proof}{\noindent\textit{Proof.}}{\hfill\(\square\)}

\title{Generalized Collatz Library}
\author{}
\date{\today}

\begin{document}
\maketitle
\section*{High-offset escape on a general base}

\begin{abstract}
A positive common-base divisor and a sufficiently divisible offset force every
nonzero signed raw orbit to escape in magnitude and eventually have the offset's
sign. Both multiplier-exponent orderings and exponent equality are
included. We prove an effective full-raw bound, reconstruct every physical
phase sweep, and show that zero is the only periodic point. The base need not
be prime, and residuals need not be units.
\end{abstract}

\section{Map and parameter range}

Throughout, division has priority in the signed raw map
\[
F(n)=
\begin{cases}
n/D,&D\mid n,\\
Qn+C,&D\nmid n,
\end{cases}
\qquad n\in\Z.
\]
The parameters satisfy
\begin{equation}\label{Parameters}
\begin{gathered}
Q=t^u,\qquad D=t^v,\qquad C=t^a s,\\
t\in\N,\quad t\ge2,\quad u,v\in\N,\quad u,v>0,\quad
a\in\N,\quad a\ge u+v,\quad s\in\Z\setminus\{0\}.
\end{gathered}
\end{equation}
Thus \(D\) is positive. There is no assumption that \(u<v\), that \(u,v\)
are coprime, or that \(t\) is prime. Nor is a unit or indivisibility
condition imposed on \(s\). A periodic point means a seed with
\(F^h(n)=n\) for some integer \(h>0\). Zero is fixed by the division branch,
even though the offset is nonzero.

The signed prime-power case
\(d=\varepsilon p^v\), \(q=p^u r\) with a high-offset phase rotation
intersects this family for \(\varepsilon=1\), \(r=1\) and both exponent
orderings. The proof below also treats composite bases and nonunit
residuals. Those stronger assertions follow from the whole-base argument,
not from silently extending prime valuation identities. Negative divisors,
nontrivial residual multipliers, and other offset regimes are not assertions
of this article.

\section{Whole-base order and legal initial divisions}

\begin{definition}\label{Order}
For \(n\ne0\), its whole-base order at \(t\) is the integer \(k\ge0\)
such that \(n=t^k w\) and \(t\nmid w\).
\end{definition}

This is not prime valuation when \(t\) is composite. For instance, \(w=2\)
is allowed at base \(t=4\): it is not divisible by four, but is not a unit
modulo four.

\begin{lemma}[Initial normalization]\label{Normalize}
Every nonzero signed seed has a finite legal division prefix ending at a
state \(x=t^k w\) with \(k<v\) and \(t\nmid w\). Its length \(r_0\) and
endpoint satisfy \(r_0\le |n|\) and \(|x|\le |n|\).
\end{lemma}
\begin{proof}
The powers \(t^j\) are unbounded: induction gives \(t^j\ge j+1\).
Consequently the nonempty set of integers \(j\ge0\) with \(t^j\mid n\)
is finite, since \(t^j\mid n\ne0\) implies \(t^j\le |n|\). Let its largest
element be \(\ell\), and write \(n=t^\ell w\). Maximality gives \(t\nmid w\).
The order is unique: if \(\ell<\ell'\) were two such orders, cancellation
would make the first residual divisible by \(t\).

Write \(\ell=r_0v+k\), \(0\le k<v\). There are exactly \(r_0\) legal initial
divisions by \(D=t^v\), ending at \(x=t^k w\). Each division has priority and
decreases the magnitude of a nonzero state. The endpoint is not divisible
by \(D\): such divisibility would imply \(t^{v-k}\mid w\), hence \(t\mid w\).
Also \(r_0\le\ell\le |n|\) and \(|x|\le |n|\).
This is the normalization proved in the public whole-base library
\cite{AscendCommonBase}.
\end{proof}

\section{The additive child}

The following lemma is stated for an arbitrary positive divisor; no
common-base presentation is required.

\begin{lemma}[Signed additive tails]\label{Additive}
Let \(d\ge2\), let \(c\ne0\) be an integer divisible by \(d\), and let
\(H\) be the raw map with coefficients \((1,c,d)\). For every signed
\(n\ne0\) there are a legal division prefix of length \(r_0\le |n|\) and
an endpoint \(x\), with \(|x|\le |n|\) and \(d\nmid x\), such that
\[
H^{r_0+j}(n)=x+jc\qquad(j\ge0).
\]
For every \(M,r\in\N\),
\[
2|n|+M+1\le r
\quad\Longrightarrow\quad
|H^r(n)|>M\ \text{ and }\ cH^r(n)>0.
\]
No nonzero signed seed has a positive return time.
\end{lemma}
\begin{proof}
Apply Lemma~\ref{Normalize} with base \(d\) and divisor exponent one.
For every \(j\ge0\), divisibility of \(x+jc\) by \(d\), together with
\(d\mid c\), would give \(d\mid x\). Thus every remaining step is the
affine branch, and induction proves the displayed tail.

If \(r\ge2|n|+M+1\), then \(j=r-r_0\ge |x|+M+1\). Integrality gives
\(c\ge1\) or \(c\le-1\). In the first case
\(x+jc\ge-|x|+j>M\); in the second case
\(x+jc\le |x|-j<-M\). Both claims follow, including the offset sign.

Suppose \(H^h(n)=n\) with \(h>0\). Iteration commutation transports this
return to the same normalized endpoint:
\(H^{r_0+h}(n)=H^{r_0}(n)=x\).
The tail formula then gives \(x+hc=x\), contradicting \(h>0\) and \(c\ne0\).
\end{proof}

\section{Exponent compression as an actual reduction}

Put
\[
m=\lfloor u/v\rfloor,\qquad b=u\bmod v,\qquad A=a-vm.
\]
Euclidean division and \eqref{Parameters} give
\begin{equation}\label{Factor}
0\le b<v,\qquad A\ge b+v,\qquad
Q=D^m t^b,\qquad C=D^m t^A s.
\end{equation}
Let \(G\) be the raw map with coefficients
\((t^b,t^A s,D)\), on the same signed integers.

\begin{lemma}[Full-time recovery]\label{Compression}
A division of \(G\) costs one raw step of \(F\), and an affine operation
costs \(m+1\). There is a positive accumulated clock \(T_k\) such that
\[
G^k(n)=F^{T_k}(n),\qquad
T_k=k+mR_k,\qquad k\le T_k\le(m+1)k,
\]
where \(R_k\) counts the affine operations among the first \(k\) steps of \(G\).
Every raw time lies in one of these blocks. If a block starts at \(x\),
each of its raw states has magnitude at least
\(\min\{|x|,|G(x)|\}\).
\end{lemma}
\begin{proof}
The guards of \(F\) and \(G\) are identical. On a division input their
outputs agree. On an affine input, \(F(x)=D^mG(x)\), followed by \(m\)
legal divisions. For \(1\le j\le m+1\), the intermediate state is
\(D^{m+1-j}G(x)\); at \(j=0\) it is \(x\).
Because \(D\ge2\), these intermediate magnitudes are at least
\(|G(x)|\), proving the stated two-endpoint lower bound.

Add these block lengths chronologically to obtain \(T_k\). Each is
positive and at most \(m+1\), so the displayed bounds hold.
For an arbitrary raw time \(r\), the increasing unbounded sequence \(T_k\)
has an index with \(T_k\le r<T_{k+1}\); its block contains \(r\).
This is the same-seed full-time recovery of the public compression theorem
\cite{AscendCompressionII}, not a one-step conjugacy.
\end{proof}

The original obligation is full-raw escape, eventual sign, periodicity,
and exact physical rays for the map in \eqref{Parameters}. Its replacement
is the same-seed compressed map in \eqref{Factor}. If \(b>0\), that map
satisfies the hypotheses of the previously proved unordered-gcd family
with \(0<b<v\) \cite{AscendRebasing}. If \(b=0\), its multiplier is one
and its nonzero offset is divisible by \(D\), so Lemma~\ref{Additive}
closes the remaining child. The cases are exhaustive.

This reduction decreases an actual resource when \(u\ge v\).
Writing \(q'=t^b\) and \(c'=t^A s\), its coefficient height satisfies
\[
\max\{|Q|,|C|\}=D^m\max\{|q'|,|c'|\}.
\]
Here \(m>0\), \(D^m>1\), and \(\max\{|q'|,|c'|\}\ge1\), so the height
strictly decreases. When \(u<v\), \(m=0\) and the map is already in the
previously proved family; no strict decrease is claimed in that case.
There is one compressed map and two solved remainder cases, not an
uncontrolled recursive family.

\section{Effective bounds at every raw time}

For \(n\ne0\), \(M\in\N\), define
\begin{equation}\label{Budget}
B(M)=
\begin{cases}
2|n|+M+1,&b=0,\\
|n|+(|n|+1+MD^{b/h})P_0,&b>0,
\end{cases}
\end{equation}
where \(h=\gcd(b,v)\) and \(P_0=b/h+v/h\) in the positive-remainder case.

\begin{lemma}[Compressed cutoff]\label{CompressedBudget}
For every \(r\ge B(M)\),
\[
|G^r(n)|>M\qquad\text{and}\qquad sG^r(n)>0.
\]
\end{lemma}
\begin{proof}
For \(b=0\), apply Lemma~\ref{Additive} to \(c'=t^A s\); this coefficient
has the same sign as \(s\). For \(b>0\), use the raw quantitative bound
from the previous gcd-rebased family \cite{AscendRebasing}, with exponents
\((b,v,A)\). Its magnitude cutoff is exactly the second line of
\eqref{Budget}. Its sign cutoff is
\(|n|+(|n|+1)P_0\), which is at most \(B(M)\). Thus the same seed, map,
and raw clock satisfy both claims.
\end{proof}

\begin{theorem}[Full-raw escape and sign]\label{Escape}
Under \eqref{Parameters}, for every signed \(n\ne0\), every \(M\in\N\),
and every raw time \(r\),
\[
B(M)(m+1)\le r\quad\Longrightarrow\quad M<|F^r(n)|.
\]
Also \(r\ge B(0)(m+1)\) implies \(sF^r(n)>0\). In particular,
\(|F^r(n)|\to\infty\) along the full raw sequence, and the offset sign
holds at all sufficiently late raw times.
\end{theorem}
\begin{proof}
Set \(K=B(M)\). Lemma~\ref{Compression} gives \(T_K\le(m+1)K\le r\).
Choose the block containing \(r\); its compressed index \(k\) must satisfy
\(k\ge K\), since \(T_{k+1}\le T_K\) would otherwise contradict
\(r<T_{k+1}\). Lemma~\ref{CompressedBudget} makes both
\(|G^k(n)|\) and \(|G^{k+1}(n)|\) exceed \(M\).
The two-endpoint lower bound in Lemma~\ref{Compression} therefore proves
the magnitude estimate at this arbitrary raw time. The explicit cutoff
is finite for each \(M\), proving the full-sequence limit.

For the sign, sample \(G^{B(0)}(n)\) at physical time \(T_{B(0)}\).
It has the sign of \(s\). Every subsequent raw division preserves that
sign because \(D>0\), and every affine operation preserves it because
\(Q>0\) and \(C\) has the same sign as \(s\).
Hence the sign persists at every later time, including every
\(r\ge(m+1)B(0)\).
\end{proof}

\begin{theorem}[Periodic points]\label{Periodic}
Under \eqref{Parameters}, a signed integer is periodic if and only if it is zero.
\end{theorem}
\begin{proof}
Zero is fixed by division priority. If \(F^h(n)=n\) with \(h>0\) and \(n\ne0\),
then iteration gives \(F^{jh}(n)=n\) for every \(j\ge0\).
Choose \(j\) so that \(jh\) exceeds the finite cutoff in
Theorem~\ref{Escape} with \(M=|n|\). It would give \(|n|>|n|\), a contradiction.
\end{proof}

\section{All forced divisions and their physical clocks}

For a normalized state \(n=t^k w\), \(k<v\), \(t\nmid w\), define
\[
e_k=\lfloor(k+u)/v\rfloor,\qquad
E(k,n)=\bigl((k+u)\bmod v,\ (Qn+C)/D^{e_k}\bigr).
\]
This event includes one affine operation and every immediately legal division.

\begin{lemma}[Legal normalized event]\label{Event}
The event \(E(k,n)\) has raw clock \(1+e_k\), and its endpoint is
\[
t^{(k+u)\bmod v}\bigl(w+t^{a-u-k}s\bigr).
\]
It has exact whole-base order \((k+u)\bmod v\).
\end{lemma}
\begin{proof}
The input is not divisible by \(D\). Since \(a\ge u+v>u+k\),
\[
Qn+C=t^{k+u}\bigl(w+t^{a-u-k}s\bigr),
\qquad a-u-k\ge1.
\]
The bracket is not divisible by \(t\): otherwise subtracting its
\(t\)-multiple correction would make \(w\) divisible by \(t\).
Thus the affine output has exact whole-base order \(k+u\), without
requiring a prime or a unit residual. Euclidean division of \(k+u\)
gives exactly \(e_k\) legal divisions by \(t^v\), ending at the displayed
state of order below \(v\). No further division is legal.
\end{proof}

\begin{lemma}[Prefix identities]\label{Prefix}
Starting at \((k,n)\), after \(i\) normalized events the phase, residual,
and physical raw clock are respectively
\[
k_i=(k+iu)\bmod v,\qquad
w_i=w+s\sum_{j=0}^{i-1}t^{a-u-k_j},\qquad
T_i=i+\lfloor(k+iu)/v\rfloor.
\]
The endpoint is \(t^{k_i}w_i=F^{T_i}(n)\), at exact whole-base order \(k_i\).
\end{lemma}
\begin{proof}
Lemma~\ref{Event} proves the phase and residual recurrence at every endpoint,
including its admissibility for the next event. Induction gives the first
two formulas and chronological composition gives the raw endpoint identity.
For the clock, Euclidean division gives
\[
\left\lfloor\frac{k+iu}{v}\right\rfloor+
\left\lfloor\frac{(k+iu)\bmod v+u}{v}\right\rfloor
=\left\lfloor\frac{k+(i+1)u}{v}\right\rfloor .
\]
Adding one for each affine operation proves the formula by induction,
with initial clock zero because \(k<v\).
\end{proof}

\begin{lemma}[Coprime sweep]\label{Sweep}
If \(\gcd(u,v)=1\), every normalized phase \(k<v\) satisfies
\[
F^{u+v}(n)=n+\beta_k,\qquad
\beta_k=s t^{a-u-v+1+k}\sum_{j=0}^{v-1}t^j.
\]
The endpoint has the same exact whole-base order \(k\). The sweep contains
\(v\) affine edges and \(u\) division edges.
\end{lemma}
\begin{proof}
The phases \((k+ju)\bmod v\), \(0\le j<v\), are a permutation of
\(0,\ldots,v-1\). Indeed, equality of two phases gives
\(v\mid(i-j)u\); coprimality gives \(v\mid i-j\), forcing \(i=j\) in that range.
A finite injective self-map is bijective. Consequently the residual kernel is
\[
\sum_{j=0}^{v-1}t^{a-u-k_j}
=\sum_{\ell=0}^{v-1}t^{a-u-\ell}
=t^{a-u-v+1}\sum_{\ell=0}^{v-1}t^\ell.
\]
All exponents are nonnegative by \eqref{Parameters}. The final phase is \(k\).
Multiplying this residual increment by \(t^k\) gives the displayed raw increment.
Lemma~\ref{Prefix} gives clock \(v+\lfloor(k+vu)/v\rfloor=u+v\).
Thus \(v\) affine operations and \(u\) divisions are counted, not discarded.
\end{proof}

\section{Gcd rebasing and same-seed raw rays}

Put
\begin{equation}\label{Rebase}
\begin{gathered}
g=\gcd(u,v),\quad T=t^g,\quad U=u/g,\quad V=v/g,\\
A_0=\lfloor a/g\rfloor,\quad S=t^{a\bmod g}s.
\end{gathered}
\end{equation}
Here \(g,U,V>0\), \(T\ge2\), \(\gcd(U,V)=1\), and \(A_0\ge U+V\).
The last inequality follows because \(g(U+V)=u+v\le a\).
Exponent division gives
\[
Q=T^U,\qquad D=T^V,\qquad C=T^{A_0}S.
\]
These are identities for the original coefficients, not a scaling of the state
or a change in the raw guard.

\begin{theorem}[Exact physical rays]\label{Ray}
For every nonzero signed seed there exist a legal initial division prefix
\(r_0\le |n|\), an endpoint \(x\) with \(|x|\le |n|\), and a rebased phase
\(k<V\), such that, for every \(j\ge0\),
\begin{equation}\label{RayFormula}
F^{r_0+jP}(n)=x+j\beta_k,\qquad
P=U+V,\qquad
\beta_k=S T^{A_0-U-V+1+k}\sum_{\ell=0}^{V-1}T^\ell .
\end{equation}
The sampled state has exact whole-base \(T\) order \(k\) for every \(j\).
Each sweep consists of \(V\) affine edges and \(U\) division edges.
\end{theorem}
\begin{proof}
Apply Lemma~\ref{Normalize} with base \(T\) and divisor exponent \(V\).
This provides the same-seed legal prefix and the stated bounds.
Apply Lemma~\ref{Sweep} to the identical parameter presentation
\((T,U,V,A_0,S)\). The endpoint returns to the same whole-base phase,
so its next sweep is admissible. Induction and chronological iteration
give \eqref{RayFormula} for every \(j\), with the original raw time
\(r_0+j(U+V)\). No arbitrary compatible witness is required.
\end{proof}

The ray proof also permits \(s=0\), in which case the increment is zero;
the escape and periodicity theorems do not assert their conclusions there.
For \(s\ne0\), the increment in \eqref{RayFormula} is nonzero and has sign \(s\).
Its sweep clock is \(P=(u+v)/g\), not a cycle period. The full-sequence
conclusion in Theorem~\ref{Escape} was proved separately with intermediate
state coverage, rather than inferred from this sampled ray.

\section{Checks that distinguish the hypotheses}

For \((Q,C,D)=(8,32,4)\), with \((t,u,v,a,s)=(2,3,2,5,1)\),
\[
1\longmapsto40\longmapsto10\longmapsto112\longmapsto28\longmapsto7.
\]
The two normalized events have clocks two and three. The complete raw
sweep has two affine operations, three divisions, and increment six.
A one-or-two-step event would miss one of the divisions.

At exponent equality, \((Q,C,D)=(4,16,4)\) gives
\[
1\longmapsto20\longmapsto5\longmapsto36\longmapsto9,
\]
with clock two and increment four. Exponent compression gives the
additive map with coefficients \((1,4,4)\).

For \((Q,C,D)=(8,4096,4)\), the opposite-sign initial segment is
\[
-767\longmapsto-2040\longmapsto-510\longmapsto16\longmapsto4\longmapsto1.
\]
It refutes any local bound that places the initial magnitude below every
later magnitude without first requiring sign alignment. The formal
full-time proof uses late bounds at both compressed endpoints. The finite
replay uses the local intermediate norm bound only on aligned sweep starts.

The uniform offset threshold cannot be lowered by one:
\((Q,C,D)=(2,4,4)\), with \(a=u+v-1\), has the primitive raw cycle
\(1,6,16,4,1\), as already proved in the preceding library
\cite{AscendCommonBase}. Its four states are distinct, so the displayed
return has least positive period four.

\section{Formal and exact computational realizations}

The three formal modules in \texttt{QCD.CommonBase.003/lean/} are:
\begin{itemize}
\item \texttt{AdditiveDivisibleOffset.lean}: legal signed additive tails,
  effective full-raw bounds, and nonzero-return exclusion.
\item \texttt{CommonBaseGeneralHighOffset.lean}: outcomes and height decrease.
\item \texttt{CommonBaseGeneralHighOffsetRay.lean}: clocks and rays.
\end{itemize}
There are twenty-seven supporting and conclusion theorems. This is a formal
declaration count, not a count of source claims.

The exact signed integer Go replay constructs both coefficient presentations
independently and checks two complete raw sweeps, all guards, phases, affine
and division counts, and increments. Its bounded main sample uses bases
two through nine, \(1\le u,v\le8\), three exponents \(a=u+v,u+v+1,u+v+2\),
sixteen nonzero signed offset residuals, 256 ordinary signed seeds, and
forced whole-base phases. Preparation also accepts \(a\le256\); seeds are
limited to 512 bits. These execution bounds do not restrict the theorem.

The finite certificate count is independently determined as follows.
For the sixty-four exponent pairs,
\(\sum v/\gcd(u,v)=221\). Their phase-weight contributions grouped by gcd are
\[
184,\ 25,\ 4,\ 4,\ 1,\ 1,\ 1,\ 1.
\]
For a rebased whole base \(T\), exactly \(16-2\lfloor8/T\rfloor\) residuals
in \([-8,8]\setminus\{0\}\) are not divisible by \(T\).
Over original bases two through nine the sums
\(\sum\lfloor8/t^g\rfloor\) are \(12,2,1\) for \(g=1,2,3\), and zero thereafter.
Thus the forced-phase weight is
\[
8\cdot16\cdot221-2(184\cdot12+25\cdot2+4)=23764,
\]
and the complete certificate count is
\[
3\cdot16(8\cdot64\cdot256+23764)=7432128.
\]
Repeated states and repeated coefficient presentations remain separate
certificates; this is not a count of distinct orbits.

Nine tests include both exponent orderings, equality, composite nonunit
residuals, large offsets, 402-bit signed initial division prefixes,
resource refusals, multiple forced divisions, the alignment counterexample,
and the finite count. Replay uses a separate workspace per base chunk,
disjoint multiplication buffers and prepared 2048-bit storage.
For the contracted replay, \(Q<2^{32}\), \(|C|<2^{1027}\), and
\(|n|<2^{512}\). At most sixteen affine updates, with divisions decreasing
magnitude, give
\[
|z|\le Q^{16}(|n|+|C|)<2^{1540}.
\]
Phase scaling is less than \(2^{32}\), so the corresponding magnitude check
is below \(2^{1572}\), within that storage reservation. Zero warmed replay
allocations are measured on seven prepared configurations only; coefficient
setup and arbitrary-size performance are not included. Neither finite
replay nor allocation measurements replace the universal proofs.

\bibliographystyle{amsalpha}
\bibliography{CommonBaseGeneralHighOffset}

\end{document}
